\chapter{Bogoliubov Transformation and Quasiparticles}
\label{chap:bogoliubov-transformation-quasiparticles}

The previous chapter formulated quadratic pairing Hamiltonians as BdG matrices.
A BdG matrix makes the particle--hole structure transparent, but it is still not
the final solution of the many-body problem.  The final step is the
\emph{Bogoliubov transformation}: a canonical change of fermionic operators that
turns each BdG block into a collection of independent quasiparticle number operators.

This chapter gives the explicit diagonalization used throughout the later
chapters on the transverse-field XY chain, the transverse-field Ising chain, and
the Kitaev chain.  The main idea is simple.  For each independent momentum pair
$(k,-k)$, pairing mixes the empty state and the doubly occupied state.  The
Bogoliubov transformation is the fermionic rotation that chooses the correct
linear combinations of particles and holes so that the Hamiltonian becomes a collection of free quasiparticles.

\section{Independent positive-momentum sectors}
\label{sec:bt-independent-momentum-sectors}

For the spinless BdG chains considered in these notes, the momentum-space
Hamiltonian has the form
\begin{equation}
\label{eq:bt-full-momentum-hamiltonian}
    H
    =
    \sum_{k\in\mathcal{K}_\vartheta}\xi(k)c_k^\dagger c_k
    +\frac12\sum_{k\in\mathcal{K}_\vartheta}
    \left[
        \Delta_\phi(k)c_k^\dagger c_{-k}^\dagger
        +\Delta_\phi(k)^*c_{-k}c_k
    \right]
    +E_{\rm cl},
\end{equation}
where
\begin{equation}
\label{eq:bt-even-xi-odd-gap}
    \xi(-k)=\xi(k),
    \qquad
    \Delta_\phi(-k)=-\Delta_\phi(k).
\end{equation}
The allowed set $\mathcal{K}_\vartheta$ is fixed by the boundary twist
introduced in the Fourier-transform chapter.  Let $\mathcal{K}_+$ denote a set
containing exactly one representative from each non-self-conjugate pair
$\{k,-k\}$.  Self-conjugate momenta, if present, satisfy
\begin{equation}
    k\equiv -k \pmod{2\pi},
\end{equation}
and are therefore $k=0$ or $k=\pi$.

For $k\in\mathcal{K}_+$, the contribution of the pair $(k,-k)$ is
\begin{equation}
\label{eq:bt-pair-hamiltonian}
    H_k
    =
    \xi_k\left(c_k^\dagger c_k+c_{-k}^\dagger c_{-k}-1\right)
    +\Delta_k c_k^\dagger c_{-k}^\dagger
    +\Delta_k^*c_{-k}c_k,
\end{equation}
where, for compactness,
\begin{equation}
    \xi_k\equiv \xi(k),
    \qquad
    \Delta_k\equiv \Delta_\phi(k).
\end{equation}
Equivalently, with the Nambu spinor
\begin{equation}
\label{eq:bt-nambu-spinor-positive-k}
    \Psi_k
    =
    \begin{pmatrix}
        c_k\\
        c_{-k}^\dagger
    \end{pmatrix},
    \qquad
    \Psi_k^\dagger
    =
    \begin{pmatrix}
        c_k^\dagger&c_{-k}
    \end{pmatrix},
\end{equation}
one has
\begin{equation}
\label{eq:bt-pair-hamiltonian-nambu}
    H_k=\Psi_k^\dagger\mathcal{H}_{\rm BdG}(k)\Psi_k,
    \qquad
    \mathcal{H}_{\rm BdG}(k)
    =
    \begin{pmatrix}
        \xi_k&\Delta_k\\
        \Delta_k^*&-\xi_k
    \end{pmatrix}.
\end{equation}
Notice that no additional factor of $1/2$ appears in
\cref{eq:bt-pair-hamiltonian-nambu}, because $\mathcal{K}_+$ already contains
only one momentum from each pair.  This is equivalent to the all-momentum BdG
formula with a factor of $1/2$.

The positive BdG energy is
\begin{equation}
\label{eq:bt-positive-energy}
    E_k
    =
    \sqrt{\xi_k^2+\abs{\Delta_k}^2}.
\end{equation}
The goal is to find new canonical fermions $\gamma_k,\gamma_{-k}$ such that
\begin{equation}
\label{eq:bt-diagonal-goal}
    H_k
    =
    E_k\left(\gamma_k^\dagger\gamma_k+
    \gamma_{-k}^\dagger\gamma_{-k}-1\right).
\end{equation}

\section{Bogoliubov angle and coherence factors}
\label{sec:bt-coherence-factors}

Write the complex pairing amplitude as
\begin{equation}
\label{eq:bt-gap-phase-definition}
    \Delta_k=\abs{\Delta_k}\ee^{\ii\alpha_k},
    \qquad
    \alpha_k=\operatorname{Arg}\Delta_k,
\end{equation}
whenever $\Delta_k\neq0$.  Introduce the polar angle $\theta_k\in[0,\pi]$ by
\begin{equation}
\label{eq:bt-bogoliubov-angle}
    \cos\theta_k=\frac{\xi_k}{E_k},
    \qquad
    \sin\theta_k=\frac{\abs{\Delta_k}}{E_k}.
\end{equation}
The standard coherence factors are
\begin{equation}
\label{eq:bt-coherence-factors}
    u_k=\cos\frac{\theta_k}{2}
    =\sqrt{\frac{E_k+\xi_k}{2E_k}},
    \qquad
    v_k=-\ee^{\ii\alpha_k}\sin\frac{\theta_k}{2}
    =-\ee^{\ii\alpha_k}\sqrt{\frac{E_k-\xi_k}{2E_k}}.
\end{equation}
They obey
\begin{equation}
\label{eq:bt-coherence-normalization}
    \abs{u_k}^2+\abs{v_k}^2=1,
\end{equation}
and, with the above phase convention,
\begin{equation}
\label{eq:bt-coherence-ratio}
    \frac{v_k}{u_k}
    =
    -\frac{\Delta_k}{E_k+\xi_k}
\end{equation}
whenever $E_k+\xi_k\neq0$.

\begin{remark}[Gauge convention at singular points]
The expression \cref{eq:bt-coherence-factors} is a convenient local gauge.  At
points where $E_k+\xi_k=0$, the ratio in \cref{eq:bt-coherence-ratio} is not a
good coordinate even though the physical state is well defined.  One may instead
use an equivalent gauge based on $E_k-\xi_k$.  This is the same type of gauge
singularity encountered when representing a two-level eigenvector on a Bloch
sphere.
\end{remark}

\section{Canonical Bogoliubov transformation}
\label{sec:bt-canonical-transformation}

Define the quasiparticle operators by
\begin{equation}
\label{eq:bt-bogoliubov-transformation}
    \gamma_k
    =
    u_k c_k-v_k c_{-k}^\dagger,
    \qquad
    \gamma_{-k}
    =
    u_k c_{-k}+v_k c_k^\dagger.
\end{equation}
Equivalently,
\begin{equation}
\label{eq:bt-bogoliubov-transformation-matrix}
    \begin{pmatrix}
        \gamma_k\\
        \gamma_{-k}^\dagger
    \end{pmatrix}
    =
    U_k
    \begin{pmatrix}
        c_k\\
        c_{-k}^\dagger
    \end{pmatrix},
    \qquad
    U_k
    =
    \begin{pmatrix}
        u_k&-v_k\\
        v_k^*&u_k
    \end{pmatrix}.
\end{equation}
The matrix $U_k$ is unitary because of
\cref{eq:bt-coherence-normalization}.  Therefore the transformation preserves the
canonical anticommutation relations:
\begin{equation}
\label{eq:bt-quasiparticle-car}
    \{\gamma_q,\gamma_{q'}^\dagger\}=\delta_{q,q'},
    \qquad
    \{\gamma_q,\gamma_{q'}\}=0.
\end{equation}
The inverse transformation is
\begin{equation}
\label{eq:bt-inverse-bogoliubov-transformation}
    c_k=u_k\gamma_k+v_k\gamma_{-k}^\dagger,
    \qquad
    c_{-k}=u_k\gamma_{-k}-v_k\gamma_k^\dagger.
\end{equation}

\begin{proposition}[Diagonalization of one BdG block]
\label{prop:bt-one-block-diagonalization}
With $u_k,v_k$ given by \cref{eq:bt-coherence-factors}, the Bogoliubov
transformation \cref{eq:bt-bogoliubov-transformation} diagonalizes the pair
Hamiltonian:
\begin{equation}
\label{eq:bt-pair-diagonalized}
    H_k
    =
    E_k\left(\gamma_k^\dagger\gamma_k+
    \gamma_{-k}^\dagger\gamma_{-k}-1\right).
\end{equation}
\end{proposition}

\begin{proof}
Using the Nambu notation of \cref{eq:bt-nambu-spinor-positive-k}, the
transformation is
\begin{equation}
    \Gamma_k=U_k\Psi_k,
    \qquad
    \Gamma_k=
    \begin{pmatrix}
        \gamma_k\\
        \gamma_{-k}^\dagger
    \end{pmatrix}.
\end{equation}
A direct multiplication gives
\begin{equation}
\label{eq:bt-unitary-diagonalization}
    U_k\mathcal{H}_{\rm BdG}(k)U_k^\dagger
    =
    \begin{pmatrix}
        E_k&0\\
        0&-E_k
    \end{pmatrix}
    =E_k\tau_z.
\end{equation}
Hence
\begin{align}
    H_k
    &
    =\Psi_k^\dagger\mathcal{H}_{\rm BdG}(k)\Psi_k
      =\Gamma_k^\dagger E_k\tau_z\Gamma_k\notag\\
    &
    =E_k\gamma_k^\dagger\gamma_k
     -E_k\gamma_{-k}\gamma_{-k}^\dagger
      =E_k\left(\gamma_k^\dagger\gamma_k+
      \gamma_{-k}^\dagger\gamma_{-k}-1\right),
\end{align}
where the last equality uses
$\gamma_{-k}\gamma_{-k}^\dagger=1-\gamma_{-k}^\dagger\gamma_{-k}$.
\end{proof}

Thus the BdG eigenvalues $\pm E_k$ are not two independent particle species.
They are the particle and hole branches of one canonical pair of quasiparticle
modes.  The physical excitation energies are nonnegative and are given by
$E_k$.

\section{Even-sector two-level problem}
\label{sec:bt-even-sector-two-level}

The same diagonalization can be seen directly in the many-body Hilbert space.
For $k\not\equiv -k$, the even-parity subspace is spanned by
\begin{equation}
\label{eq:bt-even-basis}
    \ket{0}_{k,-k},
    \qquad
    \ket{11}_{k,-k}=c_k^\dagger c_{-k}^\dagger\ket{0}_{k,-k}.
\end{equation}
In this basis the pair Hamiltonian is represented by
\begin{equation}
\label{eq:bt-even-sector-matrix}
    H_k^{\rm even}
    =
    \begin{pmatrix}
        -\xi_k&\Delta_k^*\\
        \Delta_k&\xi_k
    \end{pmatrix}.
\end{equation}
Its eigenvalues are $-E_k$ and $+E_k$.  The lower eigenstate is
\begin{equation}
\label{eq:bt-pair-ground-state}
    \ket{\mathrm{GS}}_k
    =
    u_k\ket{0}_{k,-k}
    +v_k c_k^\dagger c_{-k}^\dagger\ket{0}_{k,-k},
\end{equation}
with the same $u_k,v_k$ as in \cref{eq:bt-coherence-factors}.  Indeed,
\begin{equation}
    \frac{v_k}{u_k}=-\frac{\Delta_k}{E_k+\xi_k}
\end{equation}
is exactly the lower-eigenvector equation for
\cref{eq:bt-even-sector-matrix}.

The odd-parity states
\begin{equation}
    c_k^\dagger\ket{0}_{k,-k},
    \qquad
    c_{-k}^\dagger\ket{0}_{k,-k}
\end{equation}
are not mixed by the pairing term.  Relative to the ground state of the pair,
they are one-quasiparticle excitations with energy cost $E_k$.  The doubly
excited state $\gamma_k^\dagger\gamma_{-k}^\dagger\ket{\mathrm{GS}}_k$ has energy
cost $2E_k$.

\section{BCS product ground state}
\label{sec:bt-bcs-product-ground-state}

The full quasiparticle vacuum is defined by
\begin{equation}
\label{eq:bt-quasiparticle-vacuum-definition}
    \gamma_q\ket{\mathrm{GS}}=0
    \qquad
    \text{for all quasiparticle modes }q.
\end{equation}
For the non-self-conjugate momentum pairs, it is the product of the pair ground
states:
\begin{equation}
\label{eq:bt-bcs-product-state}
    \ket{\mathrm{GS}}
    =
    \prod_{k\in\mathcal{K}_+}
    \left(
        u_k+v_k c_k^\dagger c_{-k}^\dagger
    \right)
    \ket{0},
\end{equation}
up to the separate treatment of self-conjugate modes.  The ordering of the
product over $k$ must be fixed once and for all, since fermion creation operators
anticommute.  A different ordering changes only an overall sign convention for
the many-body state.

The diagonal Hamiltonian is
\begin{equation}
\label{eq:bt-full-diagonal-hamiltonian}
    H
    =
    E_{\rm vac}
    +
    \sum_{k\in\mathcal{K}_+}E_k
    \left(
        \gamma_k^\dagger\gamma_k+
        \gamma_{-k}^\dagger\gamma_{-k}
    \right)
    +H_{\rm self},
\end{equation}
where
\begin{equation}
\label{eq:bt-paired-sector-vacuum-energy}
    E_{\rm vac}=E_{\rm cl}-\sum_{k\in\mathcal{K}_+}E_k
\end{equation}
modulo the constant convention used in the original Hamiltonian.  The term
$H_{\rm self}$ denotes possible $k=0$ or $k=\pi$ contributions, discussed below.

\begin{principle}[Meaning of the BCS product]
The BCS ground state is not a fixed-particle-number state.  It is a coherent
superposition of states with different fermion numbers but the same fermion
parity.  This is natural for BdG Hamiltonians, which conserve fermion parity
rather than fermion number.
\end{principle}

For each pair $(k,-k)$, the occupation and anomalous expectation values are
\begin{equation}
\label{eq:bt-bcs-basic-contractions}
    \langle c_k^\dagger c_k\rangle=\abs{v_k}^2,
    \qquad
    \langle c_{-k}^\dagger c_{-k}\rangle=\abs{v_k}^2,
    \qquad
    \langle c_{-k}c_k\rangle=u_kv_k.
\end{equation}
Using \cref{eq:bt-coherence-factors}, these become
\begin{equation}
\label{eq:bt-bcs-basic-contractions-explicit}
    \abs{v_k}^2=\frac12\left(1-\frac{\xi_k}{E_k}\right),
    \qquad
    u_kv_k=-\frac{\Delta_k}{2E_k}.
\end{equation}
These contractions are the input for the covariance-matrix and Pfaffian methods
developed in the next chapter.

\section{Self-conjugate momenta and parity bookkeeping}
\label{sec:bt-self-conjugate-momenta}

If the allowed momentum set contains $q=0$ or $q=\pi$, then $q=-q$ modulo
$2\pi$.  For spinless odd-parity pairing,
\begin{equation}
    \Delta_q=0,
\end{equation}
so such a mode cannot be paired with itself.  Its BdG contribution is instead an
ordinary one-mode problem.  In a symmetric BdG convention one writes
\begin{equation}
\label{eq:bt-self-conjugate-bdg-term}
    H_q=\xi_q\left(c_q^\dagger c_q-\frac12\right).
\end{equation}
Equivalently,
\begin{equation}
\label{eq:bt-self-conjugate-diagonal}
    H_q=\abs{\xi_q}\left(\eta_q^\dagger\eta_q-\frac12\right),
\end{equation}
where
\begin{equation}
    \eta_q=
    \begin{cases}
        c_q, & \xi_q>0,\\
        c_q^\dagger, & \xi_q<0.
    \end{cases}
\end{equation}
Thus the ground-state occupation of the self-conjugate mode is
\begin{equation}
\label{eq:bt-self-conjugate-occupation}
    n_q^{\rm GS}
    =
    \begin{cases}
        0, & \xi_q>0,\\
        1, & \xi_q<0.
    \end{cases}
\end{equation}
These modes are negligible for most thermodynamic bulk quantities but essential
for finite-size spectra and fermion-parity selection.

In Jordan--Wigner images of spin chains, the physical Hilbert space splits into
fermion-parity sectors.  The parity sector fixes the boundary condition, and the
boundary condition fixes whether self-conjugate momenta are present.  Therefore a
careful finite-size solution requires the following sequence:
\begin{equation}
\label{eq:bt-finite-size-sequence}
\begin{aligned}
    \text{spin parity sector}
    &\longrightarrow
    \text{fermion boundary condition}
    \longrightarrow
    \text{allowed momentum set}\\
    &\longrightarrow
    \text{Bogoliubov diagonalization}.
\end{aligned}
\end{equation}
Only after this sequence can one compare the vacuum energies of different parity
sectors to determine the true finite-size ground state.

\section{Role of the Fourier-gauge phase}
\label{sec:bt-fourier-gauge-phase}

The Fourier convention used in these notes allows
\begin{equation}
\label{eq:bt-fourier-phase-convention}
    c_k^{(\phi)}=\ee^{-\ii\phi}c_k^{(0)},
    \qquad
    c_j=\frac{\ee^{\ii\phi}}{\sqrt{L}}
    \sum_k \ee^{\ii kj}c_k^{(\phi)}.
\end{equation}
As explained in the previous chapter, this is a Nambu gauge choice.  It changes
\begin{equation}
\label{eq:bt-gap-phase-under-fourier-gauge}
    \Delta_k^{(\phi)}=\ee^{-2\ii\phi}\Delta_k^{(0)},
\end{equation}
but leaves the spectrum invariant:
\begin{equation}
    E_k=\sqrt{\xi_k^2+\abs{\Delta_k^{(\phi)}}^2}.
\end{equation}
Since the coherence factor $v_k$ carries the phase of $\Delta_k$, it transforms
as
\begin{equation}
\label{eq:bt-v-gauge-transformation}
    v_k^{(\phi)}=\ee^{-2\ii\phi}v_k^{(0)},
    \qquad
    u_k^{(\phi)}=u_k^{(0)},
\end{equation}
within the convention \cref{eq:bt-coherence-factors}.  The BCS factor transforms
covariantly:
\begin{equation}
\label{eq:bt-bcs-factor-gauge-invariance}
    u_k^{(\phi)}+v_k^{(\phi)}
    c_k^{(\phi)\dagger}c_{-k}^{(\phi)\dagger}
    =
    u_k^{(0)}+v_k^{(0)}c_k^{(0)\dagger}c_{-k}^{(0)\dagger}.
\end{equation}
Thus the many-body state is unchanged; only its coordinate representation in
Nambu space has been rotated.

This observation is especially useful for a Kitaev chain with complex pairing,
\begin{equation}
\label{eq:bt-complex-kitaev-gap}
    \Delta=\abs{\Delta}\ee^{\ii\theta_\Delta}.
\end{equation}
The gap function derived in the previous chapter is
\begin{equation}
\label{eq:bt-kitaev-gap-with-phase}
    \Delta_\phi(k)
    =
    2\ii\abs{\Delta}\ee^{\ii(\theta_\Delta-2\phi)}\sin k.
\end{equation}
Choosing
\begin{equation}
\label{eq:bt-kitaev-phase-gauge-choice}
    \phi=\frac{\theta_\Delta}{2}
\end{equation}
sets
\begin{equation}
    \Delta_\phi(k)=2\ii\abs{\Delta}\sin k.
\end{equation}
For a single uniform Kitaev chain, this removes the uniform superconducting phase
from the BdG block.  The phase becomes physical only when it is relative to
another phase, spatially dependent, or tied to a boundary twist or flux that
cannot be removed by a single global gauge choice.

\section{Application: quasiparticles of the Kitaev chain}
\label{sec:bt-kitaev-quasiparticles}

For the Kitaev chain in the convention of the previous chapter,
\begin{equation}
\label{eq:bt-kitaev-xi-delta}
    \xi_k=-\mu-2t\cos k,
    \qquad
    \Delta_k=2\ii\abs{\Delta}\ee^{\ii(\theta_\Delta-2\phi)}\sin k.
\end{equation}
The quasiparticle dispersion is
\begin{equation}
\label{eq:bt-kitaev-quasiparticle-spectrum}
    E_k=
    \sqrt{(\mu+2t\cos k)^2+4\abs{\Delta}^2\sin^2 k}.
\end{equation}
The Bogoliubov angle is determined by
\begin{equation}
\label{eq:bt-kitaev-angle}
    \cos\theta_k=\frac{-\mu-2t\cos k}{E_k},
    \qquad
    \sin\theta_k=\frac{2\abs{\Delta}\abs{\sin k}}{E_k}.
\end{equation}
The phase entering $v_k$ is
\begin{equation}
\label{eq:bt-kitaev-alpha}
    \alpha_k
    =
    \operatorname{Arg}\left[
        2\ii\abs{\Delta}\ee^{\ii(\theta_\Delta-2\phi)}\sin k
    \right].
\end{equation}
In the gauge $\phi=\theta_\Delta/2$, this reduces to the phase of
$\ii\sin k$.  Hence the only remaining phase discontinuity is the sign change of
$\sin k$, which reflects the odd-parity nature of spinless $p$-wave pairing.

The gap closes when both terms inside \cref{eq:bt-kitaev-quasiparticle-spectrum}
vanish:
\begin{equation}
    \sin k=0,
    \qquad
    \mu+2t\cos k=0.
\end{equation}
Thus, for real nonzero $t$, the bulk phase boundaries are
\begin{equation}
    \mu=-2t
    \quad(k=0),
    \qquad
    \mu=2t
    \quad(k=\pi).
\end{equation}
The Bogoliubov quasiparticles become gapless at these transition points.

\section{Application: quasiparticles of the XY and TFIM chains}
\label{sec:bt-xy-tfim-quasiparticles}

For the transverse-field XY chain in the normalization fixed earlier, the BdG
vector can be chosen as
\begin{equation}
\label{eq:bt-xy-d-vector}
    \bm{d}_{\rm XY}(k)
    =
    \bigl(0,-2\gamma\sin k,2(h-
    \cos k)\bigr).
\end{equation}
Therefore
\begin{equation}
\label{eq:bt-xy-xi-gap}
    \xi_k=2(h-\cos k),
    \qquad
    \Delta_k=2\ii\gamma\sin k,
\end{equation}
up to the same Nambu-gauge conventions discussed above.  The quasiparticle
energy is
\begin{equation}
\label{eq:bt-xy-quasiparticle-spectrum}
    E_k=2\sqrt{(h-\cos k)^2+\gamma^2\sin^2 k}.
\end{equation}
The transverse-field Ising chain is obtained by setting $\gamma=1$:
\begin{equation}
\label{eq:bt-tfim-spectrum}
    E_k^{\rm TFIM}
    =2\sqrt{(h-\cos k)^2+\sin^2 k}
    =2\sqrt{1+h^2-2h\cos k}.
\end{equation}
The gap closes at $h=1$ and $k=0$ in the thermodynamic limit.  With the
normalization used here, the excitation gap behaves as
\begin{equation}
\label{eq:bt-tfim-gap}
    \Delta_{\rm gap}=2\abs{1-h}
\end{equation}
near the critical point.

The XX chain corresponds to $\gamma=0$.  In that case the pairing vanishes,
\begin{equation}
    \Delta_k=0,
\end{equation}
and the Bogoliubov transformation reduces to an ordinary choice of occupied and
empty fermion modes.  Thus the XX chain is a number-conserving free-fermion
problem, while the TFIM and generic XY chain are genuine BdG pairing problems.

\section{Summary}
\label{sec:bt-summary}

Each non-self-conjugate momentum pair $(k,-k)$ of a spinless BdG chain is solved
by a canonical Bogoliubov transformation.  Starting from
\begin{equation}
    H_k
    =
    \xi_k\left(n_k+n_{-k}-1\right)
    +\Delta_k c_k^\dagger c_{-k}^\dagger
    +\Delta_k^*c_{-k}c_k,
\end{equation}
one defines
\begin{equation}
    E_k=\sqrt{\xi_k^2+\abs{\Delta_k}^2},
    \qquad
    u_k=\sqrt{\frac{E_k+\xi_k}{2E_k}},
    \qquad
    v_k=-\ee^{\ii\operatorname{Arg}\Delta_k}
    \sqrt{\frac{E_k-\xi_k}{2E_k}}.
\end{equation}
The quasiparticles are
\begin{equation}
    \gamma_k=u_kc_k-v_kc_{-k}^\dagger,
    \qquad
    \gamma_{-k}=u_kc_{-k}+v_kc_k^\dagger,
\end{equation}
and the Hamiltonian becomes
\begin{equation}
    H_k=E_k\left(\gamma_k^\dagger\gamma_k+
    \gamma_{-k}^\dagger\gamma_{-k}-1\right).
\end{equation}
The quasiparticle vacuum is the BCS product state
\begin{equation}
    \ket{\mathrm{GS}}
    =
    \prod_{k\in\mathcal{K}_+}
    \left(u_k+v_kc_k^\dagger c_{-k}^\dagger\right)\ket{0},
\end{equation}
with separate finite-size bookkeeping for possible $k=0$ and $k=\pi$ modes.
The phase of the pairing gap appears in $v_k$ and in the anomalous correlator,
but a uniform Fourier-gauge phase changes only the representation, not the
physical spectrum or state.  This diagonal form is the starting point for the
Gaussian correlation-function machinery developed in the next chapter.
